\documentclass[../../main.tex]{subfiles}
\begin{document}
\section{辛张量}\label{section:ch22-symplectic-tensors}

我们从线性代数开始。有限维向量空间 $V$ 上的一个 $2$-余向量 $\omega$，如果由 $\widetilde\omega(v)=v\lrcorner\,\omega$ 定义的线性映射 $\widetilde\omega:V\to V^*$ 可逆，就称它是\emph{非退化的}。

\begin{exercise}\label{exercise:ch22-01}
证明：对于有限维向量空间 $V$ 上的 $2$-余向量 $\omega$，下列条件等价：
\begin{enumerate}[(a)]
\item $\omega$ 是非退化的。
\item 对每个非零 $v\in V$，存在 $w\in V$ 使得 $\omega(v,w)\neq0$。
\item 借助某一组（因而任意一组）基表示 $\omega$ 的矩阵 $(\omega_{ij})$ 是非奇异的。
\end{enumerate}
\end{exercise}

非退化的 $2$-余向量称为\emph{辛张量}。赋予了一个特定辛张量的向量空间 $V$ 称为\emph{辛向量空间}。（辛张量也常被称为“辛形式”，因为它特别是一个双线性型。但为避免混淆，我们把这个名称留给下面将要定义的略有不同的东西。）

\begin{example}\label{example:ch22-standard-symplectic-tensor}
设 $V$ 是维数为 $2n$ 的向量空间，取一组基记作 $(A_1,B_1,\dots,A_n,B_n)$。令 $(\alpha_1,\beta_1,\dots,\alpha_n,\beta_n)$ 表示 $V^*$ 中相应的对偶基，并令 $\omega\in\Lambda^2(V^*)$ 为下式定义的 $2$-余向量
\begin{align}\label{align:ch22-standard-symplectic-tensor}
\omega=\sum_{i=1}^n \alpha_i\wedge\beta_i.
\end{align}
注意 $\omega$ 在基向量上的作用由
\begin{align}
\omega(A_i,A_j)=\omega(B_i,B_j)=0,\qquad
\omega(A_i,B_j)=-\omega(B_j,A_i)=\delta_{ij}
\label{align:ch22-standard-symplectic-tensor-2}
\end{align}
给出。若 $v=a_iA_i+b_iB_i\in V$ 满足对所有 $w\in V$ 都有 $\omega(v,w)=0$，则 $0=\omega(v,B_i)=a_i$ 且 $0=\omega(v,A_i)=-b_i$，从而 $v=0$。因此 $\omega$ 是非退化的，因而是一个辛张量。
\end{example}

考察 $\dim V=2$ 这一特殊情形是有益的。在此情形中，沿用前一例的记号，$\alpha_1\wedge\beta_1$ 是非退化的，并且每个 $2$-余向量都是 $\alpha_1\wedge\beta_1$ 的倍数。因此二维向量空间上每个非零 $2$-余向量都是辛的。

若 $(V,\omega)$ 是辛向量空间而 $S\subset V$ 是任意线性子空间，我们把 $S$ 的\emph{辛补} $S^\omega$ 定义为子空间
\begin{align*}
S^\omega=\{v\in V:\omega(v,w)=0\text{ 对所有 }w\in S\}.
\end{align*}
正如记号所提示的，辛补类似于内积空间中的正交补。与内积情形一样，$\dim S^\omega$ 等于 $S$ 的余维数，如下一个引理所表明。

\begin{lemma}\label{lemma:ch22-symplectic-complement-dim}
设 $(V,\omega)$ 是辛向量空间。对任意线性子空间 $S\subset V$，有 $\dim S+\dim S^\omega=\dim V$。
\end{lemma}

\begin{proof}
设 $S\subset V$ 是子空间，定义线性映射 $\Phi:V\to S^*$ 为 $\Phi(v)=(v\lrcorner\,\omega)|_S$，或等价地
\begin{align*}
\Phi(v)(w)=\omega(v,w),\qquad v\in V,\ w\in S.
\end{align*}
设 $\varphi$ 是 $S^*$ 的任意元素，并令 $\widetilde\varphi\in V^*$ 是 $\varphi$ 在整个 $V$ 上的任意延拓。由于由 $v\mapsto v\lrcorner\,\omega$ 定义的映射 $\widetilde\omega:V\to V^*$ 是同构，故存在 $v\in V$ 使得 $v\lrcorner\,\omega=\widetilde\varphi$。于是 $\Phi(v)=\varphi$，因而 $\Phi$ 是满射。由秩--零维定律，$S^\omega=\ker\Phi$ 的维数等于 $\dim V-\dim S^*=\dim V-\dim S$。
\end{proof}

\begin{exercise}\label{exercise:ch22-04}
设 $(V,\omega)$ 是辛向量空间，$S\subset V$ 是线性子空间。证明 $(S^\omega)^\omega=S$。
\end{exercise}

辛补在一个重要方面不同于正交补：虽然在内积空间中总有 $S\cap S^\perp=\{0\}$，但在辛向量空间中这未必成立。事实上，若 $S$ 是一维的，则 $\omega$ 是交错的这一事实迫使对每个 $v\in S$ 都有 $\omega(v,v)=0$，故 $S\subset S^\omega$。把这一想法再推进一步，线性子空间 $S\subset V$ 称为
\begin{itemize}
\item \emph{辛的}，若 $S\cap S^\omega=\{0\}$；
\item \emph{迷向的}，若 $S\subset S^\omega$；
\item \emph{余迷向的}，若 $S\supset S^\omega$；
\item \emph{Lagrangian 的}，若 $S=S^\omega$。
\end{itemize}

\begin{proposition}\label{proposition:ch22-subspace-types}
设 $(V,\omega)$ 是辛向量空间，$S\subset V$ 是线性子空间。
\begin{enumerate}[(a)]
\item $S$ 是辛的当且仅当 $S^\omega$ 是辛的。
\item $S$ 是辛的当且仅当 $\omega|_S$ 是非退化的。
\item $S$ 是迷向的当且仅当 $\omega|_S=0$。
\item $S$ 是余迷向的当且仅当 $S^\omega$ 是迷向的。
\item $S$ 是 Lagrang 的当且仅当 $\omega|_S=0$ 且 $\dim S=\tfrac12\dim V$。
\end{enumerate}
\end{proposition}

\begin{proof}
习题 22-1。
\end{proof}

\begin{exercise}\label{exercise:ch22-06}
设 $(V,\omega)$ 是\cref{example:ch22-standard-symplectic-tensor}中描述的维数为 $2n$ 的辛向量空间，并设 $k$ 是满足 $0\leqslant k\leqslant n$ 的整数。
\begin{enumerate}[(a)]
\item 证明 $\operatorname{span}(A_1,B_1,\dots,A_k,B_k)$ 是辛的。
\item 证明 $\operatorname{span}(A_1,\dots,A_k)$ 是迷向的。
\item 证明 $\operatorname{span}(A_1,\dots,A_n,B_1,\dots,B_k)$ 是余迷向的。
\item 证明 $\operatorname{span}(A_1,\dots,A_n)$ 是 Lagrang 的。
\item 若 $n\geqslant3$，$\operatorname{span}(A_1,A_2,B_1)$ 属于四类子空间中的哪一类（如果有的话）？
\end{enumerate}
\end{exercise}

\cref{example:ch22-standard-symplectic-tensor}中定义的辛张量 $\omega$，结果证明是所有辛张量的原型，如下一个命题所表明。这可以看作 Gram--Schmidt 算法的一个辛版本。

\begin{proposition}[辛张量的典范形式]\label{proposition:ch22-symplectic-canonical-form}
设 $\omega$ 是 $m$ 维向量空间 $V$ 上的辛张量。则 $V$ 的维数 $m$ 为偶数 $m=2n$，并且存在 $V$ 的一组基，使 $\omega$ 具有式\eqref{align:ch22-standard-symplectic-tensor}的形式。
\end{proposition}

\begin{proof}
张量 $\omega$ 在基 $(A_1,B_1,\dots,A_n,B_n)$ 下具有式\eqref{align:ch22-standard-symplectic-tensor}的形式，当且仅当它在基向量上的作用由式\eqref{align:ch22-standard-symplectic-tensor-2}给出。我们对 $m=\dim V$ 作归纳来证明：存在具有该性质的基。

$m=0$ 时无需证明。设 $(V,\omega)$ 是维数 $m\geqslant1$ 的辛向量空间，并假设命题对维数小于 $m$ 的所有辛向量空间成立。任取 $V$ 中一个非零向量 $A_1$。由于 $\omega$ 非退化，存在 $B_1\in V$ 使 $\omega(A_1,B_1)\neq0$。必要时用一个常数乘 $B_1$，可设 $\omega(A_1,B_1)=1$。因为 $\omega$ 是交错的，$B_1$ 不可能是 $A_1$ 的倍数，故集合 $\{A_1,B_1\}$ 线性无关，从而 $\dim V\geqslant2$。

令 $S\subset V$ 为 $\{A_1,B_1\}$ 张成的子空间。由\cref{lemma:ch22-symplectic-complement-dim}，$\dim S^\omega=m-2$。由于 $\omega|_S$ 非退化，由\cref{proposition:ch22-subspace-types}可知 $S$ 是辛的，因而 $S^\omega$ 也是辛的。由归纳假设，$S^\omega$ 是偶维的，且存在 $S^\omega$ 的一组基 $(A_2,B_2,\dots,A_n,B_n)$ 使对 $2\leqslant i,j\leqslant n$ 满足相应的作用关系。容易推得 $(A_1,B_1,A_2,B_2,\dots,A_n,B_n)$ 就是 $V$ 所需的基。
\end{proof}

由于这个命题，若 $(V,\omega)$ 是辛向量空间，$V$ 的一组基 $(A_1,B_1,\dots,A_n,B_n)$ 在满足上述作用关系时称为\emph{辛基}，它等价于 $\omega$ 在对偶基下由式\eqref{align:ch22-standard-symplectic-tensor}给出。该命题于是表明每个辛向量空间都有一组辛基。

这给出了 $2$-余向量非退化的另一个有用判别法。对交错张量 $\omega$，记号 $\omega^k$ 表示 $k$ 重外积 $\omega\wedge\cdots\wedge\omega$。

\begin{proposition}\label{proposition:ch22-omega-nonzero}
设 $V$ 是 $2n$ 维向量空间，$\omega\in\Lambda^2(V^*)$。则 $\omega$ 是辛张量当且仅当 $\omega^n\neq0$。
\end{proposition}

\begin{proof}
先设 $\omega$ 是辛张量。令 $(A_i,B_i)$ 为 $V$ 的一组辛基，并在对偶余标架下写 $\omega=\sum_i\alpha_i\wedge\beta_i$。于是 $\omega^n=\sum_I\alpha_{i_1}\wedge\beta_{i_1}\wedge\cdots\wedge\alpha_{i_n}\wedge\beta_{i_n}$，其中 $I=(i_1,\dots,i_n)$ 遍历所有长度为 $n$ 的多重指标。该和中凡 $I$ 有重复指标的项都为零，因为 $\alpha_i\wedge\alpha_i=0$。存活下来的项恰是那些 $I$ 为 $(1,\dots,n)$ 的排列的项，而这些项彼此相等，因为 $2$-形式在外积下彼此交换。因此
\begin{align*}
\omega^n=n!\big(\alpha_1\wedge\beta_1\wedge\cdots\wedge\alpha_n\wedge\beta_n\big)\neq0.
\end{align*}

反之，设 $\omega$ 是退化的。则存在非零向量 $v\in V$ 使 $v\lrcorner\,\omega=\widetilde\omega(v)=0$。由于用 $v$ 作内乘是反导子，这意味着 $v\lrcorner\,(\omega^n)=n(v\lrcorner\,\omega)\wedge\omega^{n-1}=0$。我们可把 $v$ 扩充为 $V$ 的一组基 $(E_1,E_2,\dots,E_{2n})$ 并令 $E_1=v$，于是 $\omega^n(E_1,\dots,E_{2n})=0$，从而 $\omega^n=0$。
\end{proof}
\end{document}
